% Generated draft: requires mathematical and source review.
\documentclass[11pt]{article}
\usepackage[margin=25mm]{geometry}
\usepackage{fontspec}
\setmainfont{Latin Modern Roman}
\usepackage{amsmath,amssymb,mathtools,unicode-math}
\setmathfont{Latin Modern Math}
\usepackage{enumitem,needspace,xurl,hyperref}
\setlength{\parindent}{0pt}
\setlength{\parskip}{5pt}
\setlength{\emergencystretch}{3em}
\newcommand{\N}{\mathbb{N}}
\newcommand{\Z}{\mathbb{Z}}
\newcommand{\Q}{\mathbb{Q}}
\newcommand{\R}{\mathbb{R}}
\newcommand{\C}{\mathbb{C}}
\newcommand{\F}{\mathbb{F}}
\newcommand{\A}{\mathbb{A}}
\newcommand{\PP}{\mathbb{P}}
\newcommand{\E}{\mathbb{E}}
\newcommand{\eps}{\varepsilon}
\newcommand{\dd}{\,\mathrm{d}}
\begin{document}
\section*{UnsolvedMath ID: 30005568}
\textbf{Source problem code:} catalogue entry
\subsection{Definitions and mathematical statement}
The supplied record does not determine a complete open problem. Its extraction joins an incomplete discussion of a converse theorem to an abelian-surface application of the Birch--Swinnerton-Dyer conjecture. The following notation and conditional fragment record the recoverable mathematical content; they do not supply the missing hypotheses or identify an open conjecture.

An abelian surface over \(\Q\) is an abelian variety of dimension \(2\). For such a surface \(A\), write \(A(\Q)\) for its group of rational points and \(\operatorname{End}_{\Q}(A)\) for its ring of endomorphisms defined over \(\Q\). Interpret the extracted notation \(X(A/\Q)\) as the Tate--Shafarevich group
\[
\mathcal{X}(A/\Q)=\ker\!\left(H^1(\Q,A)\longrightarrow\prod_v H^1(\Q_v,A)\right),
\]
where \(v\) ranges over all places of \(\Q\) and \(H^1\) denotes Galois cohomology. For a prime \(p\), its \(p\)-primary subgroup is
\[
\mathcal{X}(A/\Q)[p^\infty]=\{x\in\mathcal{X}(A/\Q):p^n x=0\text{ for some integer }n\geq 1\}.
\]
Write \(L(A,s)\) for the Hasse--Weil \(L\)-function, with \(L(A,1)\) denoting the value used in the excerpt; its analytic prerequisites are not specified there.

The excerpt labels the following as a proved result, ``Theorem 5 (L.--Zerbes)'': if \(A/\Q\) is an abelian surface satisfying \(\operatorname{End}_{\Q}(A)=\Z\), is ``deformable at \(p\)'', satisfies ``various other conditions'', and has \(L(A,1)\ne 0\), then
\[
\#A(\Q)<\infty
\qquad\text{and}\qquad
\#\mathcal{X}(A/\Q)[p^\infty]<\infty.
\]
This is only a transcription of an incompletely specified source claim. Neither deformability at \(p\) nor the additional conditions are defined in the supplied record, so this implication cannot be promoted to a fully specified theorem or conjecture.

The accompanying discussion concerns weight \((2,2)\) eigenforms associated with modular abelian surfaces and an eigenvariety parametrising \(p\)-adic families for \(\operatorname{GSp}_4\). The excerpt says these nonregular weights are not immediately covered by an unspecified preceding theorem and that the eigenvariety need not be \(\acute{e}\)tale over weight space. It reports reasons to expect failure of \(\acute{e}\)taleness at certain points associated with surfaces induced from elliptic curves over imaginary quadratic fields. Deformability at \(p\) is described as slightly weaker than \(\acute{e}\)taleness and expected to hold for imaginary-quadratic lifts; these descriptions are not definitions or established universal assertions. The source also reports then-ongoing work by James Rawson, described as the speaker's PhD student, on a computationally practical test for deformability.

The preceding, disconnected fragment attributes to Rubin a CM converse theorem in the \(p\)-ordinary case, explained there as \(p\) splitting in an imaginary quadratic field \(K\), under finiteness of the \(p\)-part of a Tate--Shafarevich group. The object and converse theorem are not identified. The subsequent words ``Recently Burungale-Tian removed'' are incomplete and specify no removed hypothesis. Consequently, no precise open assertion can be recovered from that fragment either.
\subsection{Short English statement}
The record contains an incomplete, source-attributed finiteness theorem for rational points and a Tate--Shafarevich group of an abelian surface, together with a disconnected CM converse-theorem fragment. Missing hypotheses and an undefined deformability condition prevent a precise standalone formulation. The supplied open label does not identify a recoverable open conjecture.
\subsection{Sources}
\begin{itemize}
\item Supplied local UnsolvedMath registry record, permanent ID \(30005568\), exact source code \texttt{catalogue entry}, title ``Unspecified Abelian Surface BSD Conjecture''. The registry supplies status ``open'', statement status ``unrecoverable'', and a verification note explaining that the extraction joins a proved Gross--Zagier converse discussion to an abelian-surface BSD application theorem without stating a separate open conjecture. These are supplied registry classifications, not independently verified mathematical conclusions.
\item Supplied report citation: \emph{Algebraische Zahlentheorie} (2023), Oberwolfach Reports, DOI \texttt{10.4171/owr/2023/28}, pp.\ 1543--1606. Supplied source URL and literature-source link, labelled ``Oberwolfach Report 28/2023'': \url{https://doi.org/10.4171/owr/2023/28}. The record describes Oberwolfach Reports as documenting mathematical workshops and their open-problem discussions. No precise passage location or complete theorem citation is supplied.
\item Attribution retained from the extraction: ``Theorem 5 (L.--Zerbes)'', Rubin, Burungale--Tian, and James Rawson. The initial ``L.'' is not expanded in the supplied record; no separate references for these attributions are supplied.
\item Supplied literature assessment checked on \texttt{2026-08-22}: status ``open''; the extracted fragment does not identify the conjecture, and surrounding conditional and case-specific results leave the intended incomplete claim unresolved for triage. This is supplied dated curation, not independently verified source consultation or a new literature assessment. The record itself requests independent verification of time-sensitive status and cited claims.
\end{itemize}
\end{document}
